\chapter{Hydraulique des circuits}

% ============================================================
\section{Objectifs pédagogiques}
% ============================================================

\subsection*{Compétences GCF mobilisées}
\begin{itemize}
  \item \textbf{C2-1 à C2-3} — Analyser un système hydraulique : identifier composants, chaînes d'énergie, paramètres de fonctionnement
  \item \textbf{C3-3} — Dimensionner tout ou partie d'un réseau (méthodes normées, notes de calcul justifiées)
  \item \textbf{C7-4 à C7-6} — Conduire des essais de mise en service, mesurer débits et pressions, interpréter les écarts
  \item \textbf{C8-1 à C8-4} — Vérifier les performances, identifier les dérives, concevoir des actions correctives
\end{itemize}

\subsection*{Savoirs associés mobilisés}
\begin{itemize}
  \item \textbf{A3} — Dynamique des fluides : pertes de charge, point de fonctionnement, équilibrage (niveau 3)
  \item \textbf{B4-1} — Réseaux hydrauliques : eau, eau glacée, équilibrage, sécurité, dimensionnement (niveau 3)
  \item \textbf{S5.5} — Réglementations fluides thermiques : chauffage, EFS (niveau 3)
  \item \textbf{S10.3} — Sécurités d'installation, remplissage (niveau 3)
  \item \textbf{S10.5} — Mise en service d'une installation GCF (niveau 3)
  \item \textbf{S10.6} — Critères de bon fonctionnement et optimisation (niveau 3)
\end{itemize}

% ============================================================
\section{Introduction}
% ============================================================

La maîtrise de l'hydraulique est au cœur du métier de technicien supérieur GCF. Un réseau mal dimensionné ou mal équilibré entraîne des déséquilibres de débit, une surconsommation électrique des pompes, un inconfort thermique pour les occupants et une usure prématurée des équipements.

En pratique, un réseau de chauffage ou de climatisation transporte l'énergie thermique depuis la production (chaudière, PAC, groupe froid) jusqu'aux émetteurs (radiateurs, ventilo-convecteurs, planchers chauffants) via un réseau de tuyauteries. Chaque choix — diamètre, vitesse d'eau, type de pompe, organe d'équilibrage — a des conséquences sur l'efficacité énergétique globale de l'installation.

Ce chapitre couvre les bases de la dynamique des fluides appliquées aux réseaux CVC, le calcul des pertes de charge, le dimensionnement des pompes et l'équilibrage des réseaux.

% ============================================================
\section{Rappels de mécanique des fluides}
% ============================================================

\subsection{Grandeurs fondamentales}

\begin{definition}[title=Débit volumique]
Le débit volumique $q_v$ est le volume de fluide traversant une section droite par unité de temps.
\begin{equation}
  q_v = v \cdot S \qquad \left[\si{\cubic\metre\per\second}\right]
\end{equation}
\begin{itemize}
  \item $v$ : vitesse moyenne du fluide (\si{\metre\per\second})
  \item $S$ : section droite de la conduite (\si{\square\metre})
\end{itemize}
En pratique, on utilise souvent le débit en \si{\litre\per\hour} ou \si{\cubic\metre\per\hour}.
\end{definition}

\begin{definition}[title=Débit massique]
\begin{equation}
  \dot{m} = \rho \cdot q_v \qquad \left[\si{\kilogram\per\second}\right]
\end{equation}
\begin{itemize}
  \item $\rho$ : masse volumique du fluide (\si{\kilogram\per\cubic\metre})
\end{itemize}
\end{definition}

\begin{table}[H]
  \centering
  \caption{Masse volumique et viscosité cinématique de l'eau selon la température}
  \begin{tabular}{lcccc}
    \toprule
    $T$ (°C) & $\rho$ (kg/m³) & $c_p$ (J/kg·K) & $\nu \times 10^{6}$ (m²/s) & Facteur $\rho c_p / 3600$ \\
    \midrule
    10 & 999,7 & 4192 & 1,307 & 1,164 \\
    20 & 998,2 & 4182 & 1,004 & 1,163 \\
    40 & 992,2 & 4175 & 0,658 & 1,153 \\
    60 & 983,2 & 4184 & 0,475 & 1,145 \\
    70 & 977,8 & 4190 & 0,415 & 1,139 \\
    80 & 971,8 & 4196 & 0,365 & 1,133 \\
    \bottomrule
  \end{tabular}
\end{table}

\begin{definition}[title=Pression]
La pression $P$ est la force exercée par le fluide par unité de surface. Elle s'exprime en pascal (\si{\pascal}) ou en bar (\SI{1}{bar} = \SI{100000}{\pascal}).

On distingue :
\begin{itemize}
  \item \textbf{Pression absolue} $P_{abs}$ : pression totale par rapport au vide absolu
  \item \textbf{Pression relative} $P_{rel} = P_{abs} - P_{atm}$, référencée à la pression atmosphérique
  \item \textbf{Pression différentielle} $\Delta P$ : différence de pression entre deux points du réseau
\end{itemize}
\end{definition}

\subsection{Régimes d'écoulement — Nombre de Reynolds}

\begin{definition}[title=Nombre de Reynolds]
Le nombre de Reynolds $Re$ caractérise le régime d'écoulement du fluide dans une conduite.
\begin{equation}
  Re = \frac{v \cdot D}{\nu}
\end{equation}
\begin{itemize}
  \item $v$ : vitesse du fluide (\si{\metre\per\second})
  \item $D$ : diamètre intérieur de la conduite (\si{\metre})
  \item $\nu$ : viscosité cinématique (\si{\square\metre\per\second}) — pour l'eau à 50\,°C : $\nu \approx 5{,}5 \times 10^{-7}$\,\si{\square\metre\per\second}
\end{itemize}
\end{definition}

\begin{table}[H]
  \centering
  \caption{Régimes d'écoulement selon le nombre de Reynolds}
  \begin{tabular}{lll}
    \toprule
    Régime & Condition & Caractéristiques \\
    \midrule
    Laminaire & $Re < 2300$ & Filets parallèles, peu de pertes \\
    Transitoire & $2300 < Re < 4000$ & Instable \\
    Turbulent & $Re > 4000$ & Mélange intense, pertes plus élevées \\
    \bottomrule
  \end{tabular}
\end{table}

\begin{attention}
Dans les installations CVC, l'écoulement est presque toujours \textbf{turbulent} ($Re \gg 4000$). Les formules de pertes de charge utilisées dans ce chapitre s'appliquent au régime turbulent.
\end{attention}

\subsection{Équation de Bernoulli généralisée}

L'équation de Bernoulli relie les différentes formes d'énergie du fluide entre deux points 1 et 2 d'un réseau :

\begin{equation}
  \frac{P_1}{\rho g} + \frac{v_1^2}{2g} + z_1 + H_{pompe} = \frac{P_2}{\rho g} + \frac{v_2^2}{2g} + z_2 + \Delta h_{pertes}
  \label{eq:bernoulli}
\end{equation}
\begin{itemize}
  \item $P$ : pression (\si{\pascal})
  \item $v$ : vitesse (\si{\metre\per\second})
  \item $z$ : hauteur géométrique (\si{\metre})
  \item $H_{pompe}$ : hauteur manométrique apportée par la pompe (\si{\metre})
  \item $\Delta h_{pertes}$ : pertes de charge totales (\si{\metre})
  \item $g = \SI{9{,}81}{\metre\per\square\second}$
\end{itemize}

Pour un circuit fermé bouclé (chauffage, climatisation), les termes géométriques et de pression statique se compensent. Seules les pertes de charge et la HMT de la pompe subsistent.

% ============================================================
\section{Pertes de charge}
% ============================================================

Les pertes de charge représentent la chute de pression due aux frottements du fluide dans la conduite et dans les singularités (coudes, vannes, té, etc.). Elles sont exprimées en \si{\pascal} ou en mètre de colonne d'eau (\si{mCE}).

\subsection{Pertes de charge régulières (linéaires)}

Les pertes de charge régulières sont dues aux frottements le long de la conduite. On utilise la formule de \textbf{Darcy-Weisbach} :

\begin{equation}
  \Delta P_r = \lambda \cdot \frac{L}{D} \cdot \frac{\rho v^2}{2} \qquad \left[\si{\pascal}\right]
  \label{eq:darcy}
\end{equation}
\begin{itemize}
  \item $\lambda$ : coefficient de frottement (sans dimension) — dépend de $Re$ et de la rugosité $\varepsilon/D$
  \item $L$ : longueur de la conduite (\si{\metre})
  \item $D$ : diamètre intérieur (\si{\metre})
  \item $\rho$ : masse volumique (\si{\kilogram\per\cubic\metre})
  \item $v$ : vitesse du fluide (\si{\metre\per\second})
\end{itemize}

\begin{definition}[title=Coefficient de frottement $\lambda$ — diagramme de Moody]
En régime turbulent rugueux, $\lambda$ dépend du nombre de Reynolds $Re$ et de la rugosité relative $\varepsilon/D$. Pour les tubes acier courants ($\varepsilon \approx 0{,}046$\,mm) en CVC :
\begin{itemize}
  \item DN 20 à DN 50 : $\lambda \approx 0{,}022$ à $0{,}028$
  \item DN 80 à DN 200 : $\lambda \approx 0{,}018$ à $0{,}022$
\end{itemize}
Pour les tubes en cuivre ou PER (très lisses, $\varepsilon \approx 0{,}001$\,mm) : $\lambda$ légèrement plus faible.
\end{definition}

\begin{methode}[title=Utilisation des abaques de pertes de charge]
En pratique, on utilise des abaques ou des logiciels de calcul qui donnent directement la perte de charge linéique $R$ en \si{\pascal\per\metre} en fonction du débit et du diamètre.

La perte de charge régulière est alors :
\begin{equation}
  \Delta P_r = R \cdot L \qquad \left[\si{\pascal}\right]
\end{equation}

\textbf{Règle de dimensionnement GCF :}
\begin{itemize}
  \item Vitesse recommandée : entre \SI{0{,}5}{} et \SI{1{,}5}{\metre\per\second} (pour limiter bruit et érosion)
  \item Perte de charge linéique cible : entre \SI{100}{} et \SI{300}{\pascal\per\metre}
  \item Pour les colonnes montantes d'immeuble : \SI{100}{\pascal\per\metre} maximum
  \item Pour les réseaux eau glacée (tertiaire) : \SI{200}{} à \SI{400}{\pascal\per\metre}
\end{itemize}
\end{methode}

\subsection{Pertes de charge singulières}

Les pertes de charge singulières sont dues aux perturbations locales de l'écoulement (coudes, vannes, réducteurs, tés, etc.).

\begin{equation}
  \Delta P_s = \zeta \cdot \frac{\rho v^2}{2} \qquad \left[\si{\pascal}\right]
  \label{eq:singuliere}
\end{equation}
\begin{itemize}
  \item $\zeta$ : coefficient de perte de charge singulière (adimensionnel, fourni par les fabricants)
\end{itemize}

\begin{table}[H]
  \centering
  \caption{Valeurs courantes du coefficient $\zeta$ pour les accessoires hydrauliques}
  \begin{tabular}{lll}
    \toprule
    Accessoire & Conditions & $\zeta$ typique \\
    \midrule
    Coude 90° & $R/D = 1$ & 1{,}5 \\
    Coude 90° & $R/D = 2$ & 0{,}9 \\
    Coude 45° & $R/D = 1$ & 0{,}5 \\
    Té passage direct & — & 0{,}5 \\
    Té passage en dérivation & — & 1{,}5 \\
    Vanne à soupape ouverte & — & 10 \\
    Robinet d'arrêt à boisseau sphérique & ouvert & 0{,}1 \\
    Clapet anti-retour & — & 3 à 5 \\
    Filtre en Y & propre & 2 \\
    Réducteur concentrique & — & 0{,}1 à 0{,}3 \\
    Entrée brusque de tuyauterie & — & 0{,}5 \\
    \bottomrule
  \end{tabular}
\end{table}

\subsection{Méthode de la longueur équivalente}

Une méthode simplifiée consiste à remplacer chaque singularité par une longueur équivalente de tuyau droit $L_{eq}$ qui provoquerait la même perte de charge :

\begin{equation}
  \Delta P_{total} = R \cdot (L_{réel} + \sum L_{eq})
\end{equation}

En pratique, on majore souvent la longueur développée d'un forfait de 30 à 50\,\% pour tenir compte des singularités dans les réseaux courants.

\begin{table}[H]
  \centering
  \caption{Longueurs équivalentes typiques (tube acier, eau à 70°C) — extrait abaque}
  \begin{tabular}{lccc}
    \toprule
    Accessoire & DN 20 & DN 32 & DN 50 \\
    \midrule
    Coude 90° ($R/D=1$) & 0,9 m & 1,4 m & 2,2 m \\
    Té dérivation & 0,8 m & 1,3 m & 2,1 m \\
    Robinet d'arrêt & 5,5 m & 9,0 m & 14,0 m \\
    Clapet anti-retour & 1,9 m & 3,0 m & 4,8 m \\
    \bottomrule
  \end{tabular}
\end{table}

% ============================================================
\section{Les pompes de circulation}
% ============================================================

\subsection{Caractéristiques d'une pompe}

\begin{definition}[title=Hauteur manométrique totale (HMT)]
La HMT est la pression que la pompe doit vaincre pour mettre le fluide en mouvement à travers le réseau. Elle s'exprime en mètre de colonne d'eau ou en pascal.
\begin{equation}
  HMT = \Delta h_{pertes\,totales} \qquad \left[\si{mCE}\right]
\end{equation}
Pour un circuit bouclé (chauffage, climatisation), les différences de hauteur géométrique se compensent. Seules les pertes de charge restent.
\end{definition}

\begin{definition}[title=Puissance hydraulique]
\begin{equation}
  P_{hyd} = \rho \cdot g \cdot q_v \cdot HMT \qquad \left[\si{\watt}\right]
\end{equation}
\end{definition}

\begin{definition}[title=Puissance absorbée et rendement]
\begin{equation}
  P_{abs} = \frac{P_{hyd}}{\eta_{pompe}} \qquad \left[\si{\watt}\right]
\end{equation}
Le rendement global d'une pompe en circulateur CVC est typiquement compris entre 30 et 70\,\% selon la qualité de la machine et son point de fonctionnement.
\end{definition}

\begin{definition}[title=Indice d'efficacité énergétique (IEE)]
Le règlement européen 547/2012 impose un $IEE \leq 0{,}23$ pour les circulateurs de chauffage mis sur le marché depuis 2015. Les circulateurs à vitesse variable (moteur EC) atteignent des $IEE \approx 0{,}15$ à $0{,}20$.
\end{definition}

\subsection{Courbe caractéristique de pompe}

La courbe caractéristique d'une pompe donne la HMT en fonction du débit $q_v$. Elle est fournie par le fabricant. La courbe du réseau représente la perte de charge totale en fonction du débit :

\begin{equation}
  \Delta P_{réseau} = k \cdot q_v^2
\end{equation}

\begin{definition}[title=Point de fonctionnement]
Le point de fonctionnement est l'intersection entre la courbe caractéristique de la pompe et la courbe du réseau. C'est le seul point d'équilibre stable du système.
\end{definition}

\begin{attention}
Une pompe trop puissante fonctionne à débit excessif, loin de son rendement optimal. Une pompe sous-dimensionnée ne peut pas vaincre les pertes de charge de la branche défavorisée. Le choix doit être précis, avec une légère marge de sécurité (10\,\% sur HMT et débit).
\end{attention}

\subsection{Régulation de la pompe}

\begin{table}[H]
  \centering
  \caption{Modes de régulation des circulateurs à vitesse variable}
  \begin{tabular}{p{3.5cm}p{4cm}p{4.5cm}}
    \toprule
    Mode & Principe & Usage recommandé \\
    \midrule
    Pression différentielle constante & La HMT est maintenue fixe quelle que soit la demande & Petits réseaux, radiateurs sans robinets thermostatiques \\
    Pression différentielle proportionnelle & La HMT varie proportionnellement au débit ($\Delta P \propto q_v$) & Réseaux avec robinets thermostatiques — économies jusqu'à 50\,\% \\
    Température constante & Le débit est ajusté pour maintenir $\Delta T$ & Grands réseaux, variation de charge importante \\
    \bottomrule
  \end{tabular}
\end{table}

\subsection{Association de pompes}

\begin{table}[H]
  \centering
  \caption{Association de pompes identiques}
  \begin{tabular}{lll}
    \toprule
    Montage & Effet sur le débit & Effet sur la HMT \\
    \midrule
    Série & Identique & Doublée \\
    Parallèle & Doublé & Identique \\
    \bottomrule
  \end{tabular}
\end{table}

\begin{attention}
Le doublement du débit ou de la HMT n'est valable que pour des pompes identiques sur un réseau idéal. En pratique, l'augmentation réelle est moindre à cause de la courbe de réseau. Le montage en parallèle est utilisé pour la redondance (secours) et les périodes de pointe.
\end{attention}

% ============================================================
\section{Équilibrage des réseaux hydrauliques}
% ============================================================

\subsection{Problématique du déséquilibre}

Dans un réseau ramifié, certains émetteurs reçoivent plus d'eau que prévu (les plus proches de la pompe) tandis que d'autres sont sous-alimentés (les plus éloignés). Ce déséquilibre provoque :
\begin{itemize}
  \item Une surchauffe des pièces proches de la production
  \item Un inconfort dans les pièces éloignées
  \item Une consommation électrique excessive (la pompe lutte contre un réseau mal équilibré)
\end{itemize}

\subsection{Organes d'équilibrage}

\begin{table}[H]
  \centering
  \caption{Organes d'équilibrage — caractéristiques et usages}
  \begin{tabular}{p{3.5cm}p{5cm}p{4cm}}
    \toprule
    Organe & Principe & Application \\
    \midrule
    Robinet d'équilibrage manuel & Crée une perte de charge réglable. Nécessite une mesure de débit. & Tout réseau statique \\
    Régulateur de débit automatique (PIC) & Maintient un débit constant indépendamment de la pression différentielle & Réseau à débit variable (robinets thermostatiques) \\
    Vanne de régulation pression différentielle & Maintient une $\Delta P$ constante sur un circuit & Réseaux de distribution \\
    Vanne motorisée 2 voies & Ouverture proportionnelle ou TOR commandée par régulateur & Émetteur avec régulation locale \\
    \bottomrule
  \end{tabular}
\end{table}

\begin{methode}[title=Équilibrage par la méthode de la résistance équivalente (branche index)]
\begin{enumerate}
  \item Calculer les pertes de charge de chaque branche du réseau (du départ au retour)
  \item Identifier la branche la plus défavorisée (pertes de charge les plus élevées) — c'est la branche \textbf{index}
  \item Pour chaque autre branche, calculer l'excédent de pression à absorber :
    \[
      \Delta P_{excédent} = \Delta P_{index} - \Delta P_{branche}
    \]
  \item Régler les organes d'équilibrage pour créer une perte de charge supplémentaire égale à $\Delta P_{excédent}$ sur chaque branche
  \item Vérifier les débits après réglage (débitmètre, mesure différentielle)
\end{enumerate}
\end{methode}

\subsection{Sécurités hydrauliques}

\begin{table}[H]
  \centering
  \caption{Dispositifs de sécurité des installations hydrauliques (DTU 65.11)}
  \begin{tabular}{p{3.5cm}p{4.5cm}p{4.5cm}}
    \toprule
    Dispositif & Rôle & Réglage typique \\
    \midrule
    Vase d'expansion & Absorbe la dilatation de l'eau en chauffe (eau : $+4\,\%$ de 10 à 90°C) & Pression de gonflage : $P_{stat} + 0{,}3$\,bar \\
    Soupape de sécurité & Limite la pression max en cas de surpression & 3 bar (chaudière murale) \\
    Pressostat de sécurité & Coupe le brûleur si pression anormale & Bas : 0,5 bar / Haut : 3 bar \\
    Clapet anti-retour & Empêche la circulation inverse dans les circuits en parallèle & — \\
    Purgeur automatique & Élimine l'air dissous pouvant créer des bouchons & Point haut de chaque circuit \\
    \bottomrule
  \end{tabular}
\end{table}

% ============================================================
\section{Exemples résolus}
% ============================================================

\begin{exemple}[title=Calcul des pertes de charge d'un tronçon]
\textbf{Énoncé :} Un tronçon de réseau de chauffage est constitué d'un tube acier DN 25 (diamètre intérieur \SI{27{,}2}{\milli\metre}) de longueur \SI{15}{\metre}. Il débite \SI{800}{\litre\per\hour} d'eau à 70\,°C. Ce tronçon comporte 4 coudes à 90° ($\zeta = 1{,}5$ chacun) et 1 robinet d'arrêt à boisseau ($\zeta = 0{,}1$). Masse volumique de l'eau à 70\,°C : $\rho = \SI{978}{\kilogram\per\cubic\metre}$.

À partir de l'abaque du fabricant, la perte de charge linéique pour ce diamètre et ce débit est $R = \SI{250}{\pascal\per\metre}$.

\textbf{Étape 1 : Perte de charge régulière}
\[
  \Delta P_r = R \times L = 250 \times 15 = \SI{3750}{\pascal}
\]

\textbf{Étape 2 : Vitesse du fluide}
\[
  S = \frac{\pi \times (0{,}0272)^2}{4} = \SI{5{,}81e-4}{\square\metre}
\]
\[
  q_v = \frac{800}{3600 \times 1000} = \SI{2{,}22e-4}{\cubic\metre\per\second}
\]
\[
  v = \frac{2{,}22 \times 10^{-4}}{5{,}81 \times 10^{-4}} = \SI{0{,}38}{\metre\per\second}
\]

\textbf{Étape 3 : Pression dynamique}
\[
  \frac{\rho v^2}{2} = \frac{978 \times (0{,}38)^2}{2} = \SI{70{,}6}{\pascal}
\]

\textbf{Étape 4 : Pertes de charge singulières}
\[
  \Delta P_s = \left(4 \times 1{,}5 + 1 \times 0{,}1\right) \times 70{,}6 = 6{,}1 \times 70{,}6 = \SI{430}{\pascal}
\]

\textbf{Étape 5 : Perte de charge totale}
\[
  \Delta P_{total} = 3750 + 430 = \SI{4180}{\pascal} \approx \SI{0{,}042}{\bar}
\]

\textbf{Vérification de la vitesse :} $v = 0{,}38$\,m/s — légèrement en dessous de la plage 0,5–1,5\,m/s, acceptable pour une colonne montante (critère acoustique prioritaire).
\end{exemple}

\begin{exemple}[title=Dimensionnement d'une pompe de chauffage]
\textbf{Énoncé :} Un réseau de chauffage délivre une puissance de \SI{45}{\kilo\watt} avec une température de départ à 70\,°C et retour à 55\,°C. La perte de charge totale de la boucle la plus défavorisée est \SI{18000}{\pascal}. Déterminer le débit nécessaire et choisir une pompe.

\textbf{Étape 1 : Calcul du débit}
\[
  \dot{m} = \frac{\dot{Q}}{c_p \cdot \Delta T} = \frac{45\,000}{4186 \times (70 - 55)} = \SI{0{,}717}{\kilogram\per\second}
\]
\[
  q_v = \frac{0{,}717}{978} = \SI{7{,}33e-4}{\cubic\metre\per\second} = \SI{2{,}64}{\cubic\metre\per\hour}
\]

\textbf{Étape 2 : HMT requise}
\[
  HMT = \frac{\Delta P_{réseau}}{\rho \cdot g} = \frac{18\,000}{978 \times 9{,}81} = \SI{1{,}88}{\mCE} \approx \SI{1{,}9}{\mCE}
\]

\textbf{Étape 3 : Sélection pompe (avec marge 10\,\%)}
\begin{itemize}
  \item Débit de sélection : $2{,}64 \times 1{,}10 \approx \SI{2{,}9}{\cubic\metre\per\hour}$
  \item HMT de sélection : $1{,}9 \times 1{,}10 \approx \SI{2{,}1}{\mCE}$
\end{itemize}
Rechercher dans les catalogues fabricants (Grundfos, Wilo, Armstrong) un circulateur délivrant au minimum $q_v = \SI{2{,}9}{\cubic\metre\per\hour}$ à $HMT = \SI{2{,}1}{\mCE}$.

\textbf{Résultat :} Un circulateur de type Grundfos ALPHA2 25-40 ou équivalent convient (plage : 0 à 3\,m³/h, HMT max 4\,mCE), avec régulation de vitesse intégrée pour l'efficacité énergétique ($IEE \leq 0{,}20$).
\end{exemple}

\begin{exemple}[title=Équilibrage d'un réseau bitubes — 3 radiateurs]
\textbf{Énoncé :} Un réseau bitubes alimente 3 radiateurs depuis une nourrice commune. Les pertes de charge calculées de la nourrice jusqu'à chaque radiateur sont :
\begin{itemize}
  \item Radiateur R1 (proche) : $\Delta P_{R1} = \SI{2800}{\pascal}$
  \item Radiateur R2 (intermédiaire) : $\Delta P_{R2} = \SI{5100}{\pascal}$
  \item Radiateur R3 (éloigné, branche index) : $\Delta P_{R3} = \SI{7600}{\pascal}$
\end{itemize}

\textbf{Perte de charge à créer sur chaque robinet d'équilibrage :}
\begin{itemize}
  \item R1 : $\Delta P_{rob,R1} = 7600 - 2800 = \SI{4800}{\pascal}$ → régler le robinet
  \item R2 : $\Delta P_{rob,R2} = 7600 - 5100 = \SI{2500}{\pascal}$ → régler le robinet
  \item R3 : $\Delta P_{rob,R3} = \SI{0}{\pascal}$ → robinet grand ouvert (branche index)
\end{itemize}

\textbf{La pompe doit fournir} au minimum $\SI{7600}{\pascal}$ (pertes de la branche index) plus les pertes du tronçon commun.

\textbf{Vérification :} Après réglage, les 3 radiateurs ont la même résistance hydraulique vue depuis la nourrice → même pression disponible → débits nominaux respectés.
\end{exemple}

\begin{exemple}[title=Dimensionnement d'un vase d'expansion]
\textbf{Énoncé :} Une installation de chauffage contient un volume d'eau total $V_{eau} = \SI{150}{\litre}$ (chaudière + réseau). La pression statique de remplissage est $P_{stat} = \SI{1{,}5}{\bar}$. La pression max de la soupape est $P_{soupape} = \SI{3}{\bar}$. La température passe de 10\,°C (remplissage) à 80\,°C (chauffe). Calculer le volume du vase.

\textbf{Dilatation de l'eau :} $\rho_{10°C} = 999{,}7$\,kg/m³, $\rho_{80°C} = 971{,}8$\,kg/m³.
\[
  \Delta V_{eau} = V_{eau} \times \left(\frac{\rho_{10}}{\rho_{80}} - 1\right) = 150 \times \left(\frac{999{,}7}{971{,}8} - 1\right) = 150 \times 0{,}0287 = \SI{4{,}3}{\litre}
\]

\textbf{Volume utile du vase :}
Le vase à membrane est gonflé à $P_{gonflage} = P_{stat} = \SI{1{,}5}{\bar}$.
\[
  V_{vase} = \Delta V_{eau} \times \frac{P_{soupape} + 1}{P_{soupape} - P_{stat}} = 4{,}3 \times \frac{3 + 1}{3 - 1{,}5} = 4{,}3 \times \frac{4}{1{,}5} = \SI{11{,}5}{\litre}
\]
Choisir un vase de \SI{12}{\litre} (première taille standard supérieure).
\end{exemple}

% ============================================================
\section{Schémas et diagrammes caractéristiques}
% ============================================================

\subsection{Circuit hydraulique bouclé}

\begin{figure}[H]
  \centering
  \begin{tikzpicture}[>=Stealth, font=\small, thick,
      lbl/.style={font=\footnotesize, inner sep=2pt}]

    % ---- CHAUDIÈRE (symbole ISO NF) ----
    \pic[draw=FedBlue] at (0,0) {chaudiere};
    \node[lbl, above=0.65cm] at (0,0) {Chaudière};
    \node[lbl, text=FedBlue, below=0.35cm] at (0,0) {\footnotesize condensation};

    % ---- VASE D'EXPANSION (symbole ISO) ----
    \pic[draw=FedGreen!70!black] at (2.5,1.2) {vase-expansion};
    \node[lbl, above=0.75cm] at (2.5,1.2) {Vase exp.};

    % ---- POMPE DE CIRCULATION (symbole ISO) ----
    \pic[draw=FedOrange] at (0,-2.2) {pompe};
    \node[lbl, right=0.55cm] at (0,-2.2) {Circ.};

    % ---- NOURRICE (rectangle simple — nœud de distribution) ----
    \draw[draw=FedBlue, fill=FedBlue!10, rounded corners=2pt]
      (4.5,-0.25) rectangle (5.5,0.25);
    \node[lbl] at (5.0,0) {Nourrice};

    % ---- RADIATEURS (symboles ISO NF) ----
    \pic[draw=FedRed] at (2.5,-3.5) {radiateur};
    \node[lbl, below=0.55cm] at (2.5,-3.5) {Rad. 1};

    \pic[draw=FedRed] at (5.0,-3.5) {radiateur};
    \node[lbl, below=0.55cm] at (5.0,-3.5) {Rad. 2};

    \pic[draw=FedRed] at (7.5,-3.5) {radiateur};
    \node[lbl, below=0.55cm] at (7.5,-3.5) {Rad. 3};

    % ---- VANNES QUART DE TOUR sur chaque radiateur ----
    \pic[draw=FedBlue!60] at (2.5,-2.70) {vanne-quart};
    \pic[draw=FedBlue!60] at (5.0,-2.70) {vanne-quart};
    \pic[draw=FedBlue!60] at (7.5,-2.70) {vanne-quart};

    % ============ TUYAUTERIE DÉPART (rouge) ============
    % Chaudière (est) → nourrice (ouest)
    \draw[->, draw=FedRed, very thick]
      (0.80,0) -- node[above, lbl, text=FedRed]{Départ 70\,°C} (4.5,0);
    % Vase d'expansion raccordé sur la ligne départ
    \draw[draw=FedGreen!70!black, thick] (2.5,0.70) -- (2.5,0) -- (3.5,0);
    % Nourrice → distribution verticale
    \draw[draw=FedRed, thick] (5.0,-0.25) -- (5.0,-2.35);
    \draw[draw=FedRed, thick] (5.0,-2.35) -- (2.5,-2.35) -- (7.5,-2.35);
    % Départ vers vannes
    \draw[->, draw=FedRed, thick] (2.5,-2.35) -- (2.5,-2.38);
    \draw[->, draw=FedRed, thick] (5.0,-2.35) -- (5.0,-2.38);
    \draw[->, draw=FedRed, thick] (7.5,-2.35) -- (7.5,-2.38);
    % Vannes → radiateurs
    \draw[->, draw=FedRed, thick] (2.5,-3.02) -- (2.5,-3.12);
    \draw[->, draw=FedRed, thick] (5.0,-3.02) -- (5.0,-3.12);
    \draw[->, draw=FedRed, thick] (7.5,-3.02) -- (7.5,-3.12);

    % ============ TUYAUTERIE RETOUR (bleu) ============
    \draw[draw=FedBlue, thick] (2.5,-3.88) -- (2.5,-4.6);
    \draw[draw=FedBlue, thick] (5.0,-3.88) -- (5.0,-4.6);
    \draw[draw=FedBlue, thick] (7.5,-3.88) -- (7.5,-4.6);
    \draw[draw=FedBlue, thick] (2.5,-4.6) -- (7.5,-4.6);
    \draw[draw=FedBlue, thick] (5.0,-4.6) -- (5.0,-5.1) -- (0,-5.1) -- (0,-2.62);
    \draw[->, draw=FedBlue, very thick]
      (0,-2.62) -- node[right, lbl, text=FedBlue]{Retour 55\,°C} (0,-2.62);
    \draw[->, draw=FedBlue, very thick] (0,-2.62) -- (0.0,-2.62);
    % Retour → pompe → chaudière
    \draw[draw=FedBlue, very thick] (0,-5.1) -- (0,-2.62);
    \draw[->, draw=FedBlue, very thick] (0,-1.78) -- (0,-0.42);

    % ============ LÉGENDE ============
    \begin{scope}[yshift=-6.0cm]
      \draw[FedRed,  very thick, ->] (0,0) -- (0.7,0);
      \node[right, lbl, text=FedRed]  at (0.8,0)   {Aller — eau chaude};
      \draw[FedBlue, very thick, ->]  (3.8,0) -- (4.5,0);
      \node[right, lbl, text=FedBlue] at (4.6,0)   {Retour — eau tiède};
      \node[lbl, text=FedGreen!70!black]  at (1.5,-0.55)
        {\tikz\pic[draw=FedGreen!70!black, scale=0.7]{vase-expansion}; \;Vase d'expansion};
      \node[lbl, text=FedOrange]          at (6.5,-0.55)
        {\tikz\pic[draw=FedOrange, scale=0.7]{pompe}; \;Pompe circulatrice};
    \end{scope}

  \end{tikzpicture}
  \caption{Circuit hydraulique bouclé — chauffage bitubes (symboles ISO NF GCF :
    chaudière, pompe circulatrice, vase d'expansion, radiateurs, vannes quart de tour)}
  \label{fig:circuit-boucle}
\end{figure}

\subsection{Courbe caractéristique pompe / réseau}

\begin{figure}[H]
  \centering
  \begin{tikzpicture}[>=Stealth, thick]
    % Axes
    \draw[->] (0,0) -- (7.5,0) node[right, font=\small] {$q_v$ (\si{\cubic\metre\per\hour})};
    \draw[->] (0,0) -- (0,5.5) node[above, font=\small] {$H$ (mCE)};

    % Graduations axe X
    \foreach \x/\lbl in {1/0{,}5, 2/1{,}0, 3/1{,}5, 4/2{,}0, 5/2{,}5, 6/3{,}0} {
      \draw (\x,0.05) -- (\x,-0.05);
      \node[below, font=\tiny] at (\x,0) {\lbl};
    }
    % Graduations axe Y
    \foreach \y/\lbl in {1/1, 2/2, 3/3, 4/4, 5/5} {
      \draw (0.05,\y) -- (-0.05,\y);
      \node[left, font=\tiny] at (0,\y) {\lbl};
    }

    % Courbe pompe : H = 5 - 0.2*qv^2 (parabole descendante)
    \draw[color=FedBlue, very thick, domain=0:6.5, samples=60]
      plot (\x, {5 - 0.1*\x*\x})
      node[right, font=\small, text=FedBlue] {Courbe pompe};

    % Courbe réseau : H = 0.12*qv^2 (parabole ascendante, passe par 0)
    \draw[color=FedOrange, very thick, domain=0:6.2, samples=60]
      plot (\x, {0.13*\x*\x})
      node[above, font=\small, text=FedOrange] {Courbe réseau};

    % Point de fonctionnement : intersection de 5 - 0.1*x^2 = 0.13*x^2
    % 5 = 0.23*x^2 => x^2 = 21.74 => x = 4.66, y = 0.13*21.74 = 2.83
    \draw[dashed, color=FedRed] (4.66,0) -- (4.66,2.83) -- (0,2.83);
    \filldraw[FedRed] (4.66,2.83) circle (3pt);
    \node[above right, font=\small\bfseries, text=FedRed] at (4.66,2.83) {Pt de fonctionnement};
    \node[below, font=\tiny, text=FedRed] at (4.66,0) {$q_{v,0}$};
    \node[left, font=\tiny, text=FedRed] at (0,2.83) {$H_0$};

    % Zone de bon rendement (hachures)
    \draw[color=FedGreen!60, dashed, thick] (3.5,0) -- (3.5,{5 - 0.1*3.5*3.5});
    \draw[color=FedGreen!60, dashed, thick] (5.5,0) -- (5.5,{5 - 0.1*5.5*5.5});
    \node[font=\tiny, text=FedGreen!80!black, rotate=90] at (4.5,1.0) {Zone optimale};
  \end{tikzpicture}
  \caption{Courbe caractéristique pompe et courbe de réseau — détermination du point de fonctionnement}
  \label{fig:courbe-pompe-reseau}
\end{figure}

\begin{methode}[title=Lecture du point de fonctionnement]
Le point de fonctionnement $\left(q_{v,0},\, H_0\right)$ est l'intersection des deux courbes :
\begin{itemize}
  \item \textbf{Courbe pompe} : $H_{pompe}(q_v)$ décroissante — fournie par le fabricant
  \item \textbf{Courbe réseau} : $H_{réseau} = k \cdot q_v^2$ croissante — calculée par l'installateur
\end{itemize}
Si le réseau est modifié (équilibrage, remplacement d'un émetteur), la courbe réseau change et le point de fonctionnement se déplace sur la courbe pompe.
\end{methode}

\subsection{Régimes d'écoulement — Profils de vitesse}

\begin{figure}[H]
  \centering
  \begin{tikzpicture}[>=Stealth, thick, node distance=4.5cm]

    %--- Régime laminaire (gauche) ---
    \begin{scope}[xshift=0cm]
      % Tube
      \draw[draw=FedBlue!50, fill=FedBlue!8] (0,-1.2) rectangle (4,1.2);
      \draw[very thick] (0,-1.2) -- (4,-1.2);
      \draw[very thick] (0, 1.2) -- (4, 1.2);

      % Profil parabolique : v(r) = v_max*(1-(r/R)^2), R=1.2
      % r=0.0 : vlen=1.69; r=±0.3 : vlen=1.55; r=±0.6 : vlen=1.12; r=±0.9 : vlen=0.62; r=±1.1 : vlen=0.23
      \draw[->, color=FedBlue, thick] (0.6, 0.0)  -- ++(1.69,0);
      \draw[->, color=FedBlue, thick] (0.6, 0.3)  -- ++(1.55,0);
      \draw[->, color=FedBlue, thick] (0.6,-0.3)  -- ++(1.55,0);
      \draw[->, color=FedBlue, thick] (0.6, 0.6)  -- ++(1.12,0);
      \draw[->, color=FedBlue, thick] (0.6,-0.6)  -- ++(1.12,0);
      \draw[->, color=FedBlue, thick] (0.6, 0.9)  -- ++(0.62,0);
      \draw[->, color=FedBlue, thick] (0.6,-0.9)  -- ++(0.62,0);
      \draw[->, color=FedBlue, thick] (0.6, 1.1)  -- ++(0.23,0);
      \draw[->, color=FedBlue, thick] (0.6,-1.1)  -- ++(0.23,0);

      % Légende
      \node[below, font=\small\bfseries, text=FedBlue] at (2,-1.5) {Laminaire ($Re < 2300$)};
      \node[below, font=\footnotesize] at (2,-1.9) {Profil parabolique};
      \node[right, font=\tiny] at (2.3,0) {$v_{max}$};

      % Note Reynolds
      \node[font=\footnotesize, text=FedBlue!70!black] at (2,1.6) {Filets parallèles, $\lambda = 64/Re$};
    \end{scope}

    %--- Régime turbulent (droite) ---
    \begin{scope}[xshift=5.5cm]
      % Tube
      \draw[draw=FedOrange!50, fill=FedOrange!8] (0,-1.2) rectangle (4,1.2);
      \draw[very thick] (0,-1.2) -- (4,-1.2);
      \draw[very thick] (0, 1.2) -- (4, 1.2);

      % Profil quasi-uniforme : flèches presque égales, chute rapide en paroi
      \draw[->, color=FedOrange, thick] (0.6, 0.0)  -- ++(1.69,0);
      \draw[->, color=FedOrange, thick] (0.6, 0.3)  -- ++(1.63,0);
      \draw[->, color=FedOrange, thick] (0.6,-0.3)  -- ++(1.63,0);
      \draw[->, color=FedOrange, thick] (0.6, 0.6)  -- ++(1.50,0);
      \draw[->, color=FedOrange, thick] (0.6,-0.6)  -- ++(1.50,0);
      \draw[->, color=FedOrange, thick] (0.6, 0.9)  -- ++(1.20,0);
      \draw[->, color=FedOrange, thick] (0.6,-0.9)  -- ++(1.20,0);
      \draw[->, color=FedOrange, thick] (0.6, 1.1)  -- ++(0.40,0);
      \draw[->, color=FedOrange, thick] (0.6,-1.1)  -- ++(0.40,0);

      % Perturbations schématisées
      \draw[color=FedOrange!60, thin, ->] (1.5, 0.4) -- ++(-0.1,0.25) -- ++(0.15,0.2);
      \draw[color=FedOrange!60, thin, ->] (1.5,-0.4) -- ++(0.1,-0.25) -- ++(-0.15,-0.2);
      \draw[color=FedOrange!60, thin, ->] (2.3, 0.4) -- ++(-0.1,0.25) -- ++(0.15,0.2);
      \draw[color=FedOrange!60, thin, ->] (2.3,-0.4) -- ++(0.1,-0.25) -- ++(-0.15,-0.2);

      % Légende
      \node[below, font=\small\bfseries, text=FedOrange] at (2,-1.5) {Turbulent ($Re > 4000$)};
      \node[below, font=\footnotesize] at (2,-1.9) {Profil quasi-uniforme};
      \node[font=\footnotesize, text=FedOrange!70!black] at (2,1.6) {Mélange transversal, $\lambda = f(Re,\, \varepsilon/D)$};
    \end{scope}

    % Flèche de transition au centre
    \draw[->, very thick, color=FedGreen] (4.2,0) -- node[above, font=\footnotesize, text=FedGreen]{$Re \uparrow$} (5.3,0);

  \end{tikzpicture}
  \caption{Profils de vitesse en régime laminaire et turbulent dans une conduite circulaire}
  \label{fig:profils-ecoulement}
\end{figure}

% ============================================================
\section{Éléments de réglementation}
% ============================================================

\begin{description}
  \item[DTU 65.11 (NF P52-211)] Dispositifs de sécurité des installations de chauffage central. Prescrit les dispositifs de protection (soupapes, vases d'expansion, pressostat) et leurs dimensionnements.
  \item[DTU 60.1 (NF P40-201)] Plomberie sanitaire — règles de calcul et installation. Dimensionnement des réseaux d'eau froide et chaude sanitaire.
  \item[NF EN 12828] Installations de chauffage dans les bâtiments — conception des systèmes de chauffage à eau chaude. Prescrit notamment les pressions de service, températures maximales et dispositifs de sécurité.
  \item[NF EN 14336] Mise en service des installations de chauffage à eau chaude. Procédures d'essais d'étanchéité, rinçage, remplissage et équilibrage.
  \item[NF EN 12831] Méthode de calcul des besoins en chaleur — débit à distribuer par émetteur (base du dimensionnement des tuyauteries).
  \item[RE 2020] Impose une approche de calcul de la consommation d'énergie primaire incluant les auxiliaires (pompes). Favorise les circulateurs à haut rendement (classe ErP 2015 minimum).
  \item[Règlement EU 547/2012] Définit les exigences d'écoconception des pompes à eau. Les circulateurs de chauffage doivent satisfaire l'indice d'efficacité énergétique $IEE \leq 0{,}23$.
\end{description}

\begin{tcolorbox}[
  enhanced, colback=FedBlue!3, colframe=FedBlue!50,
  fonttitle=\bfseries\sffamily, coltitle=FedBlue,
  title={FICHE MÉMO — Hydraulique des circuits}, boxrule=0.8pt, arc=2pt,
  width=\textwidth
]
\begin{multicols}{2}
\textbf{Débit volumique}
\[ q_v = v \cdot S \quad \left[\si{\cubic\metre\per\second}\right] \]
avec $v$ = vitesse [\si{\metre\per\second}], $S$ = section droite [\si{\metre\squared}]

\medskip
\textbf{Débit massique}
\[ \dot{m} = \rho \cdot q_v \quad \left[\si{\kilogram\per\second}\right] \]
avec $\rho$ = masse volumique [\si{\kilogram\per\cubic\metre}]

\medskip
\textbf{Nombre de Reynolds}
\[ Re = \frac{v \cdot D}{\nu} \]
$Re < 2300$ : laminaire \quad $Re > 4000$ : turbulent \\
avec $\nu$ = viscosité cinématique [\si{\metre\squared\per\second}]

\medskip
\textbf{Pertes de charge régulières (Darcy-Weisbach)}
\[ \Delta P_r = \lambda \cdot \frac{L}{D} \cdot \frac{\rho\, v^2}{2} \quad \left[\si{\pascal}\right] \]
avec $\lambda$ = coeff. de frottement, $L$ = longueur [\si{\metre}]

\columnbreak

\textbf{Pertes de charge singulières}
\[ \Delta P_s = \zeta \cdot \frac{\rho\, v^2}{2} \quad \left[\si{\pascal}\right] \]
avec $\zeta$ = coefficient de singularité (coude 90° : $\zeta \approx 1{,}5$)

\medskip
\textbf{Méthode de la longueur équivalente}
\[ \Delta P_{total} = R \cdot \bigl(L_{réel} + \textstyle\sum L_{eq}\bigr) \]
avec $R$ = perte de charge linéique [\si{\pascal\per\metre}] \\
Règle GCF : cible $R$ entre \SI{100}{} et \SI{300}{\pascal\per\metre}

\medskip
\textbf{HMT — Hauteur manométrique totale}
\[ HMT = \Delta h_{pertes\,totales} \quad \left[\si{mCE}\right] \]

\medskip
\textbf{Puissance hydraulique}
\[ P_{hyd} = \rho \cdot g \cdot q_v \cdot HMT \quad \left[\si{\watt}\right] \]

\medskip
\textbf{Puissance absorbée par la pompe}
\[ P_{abs} = \frac{P_{hyd}}{\eta_{pompe}} \quad \left[\si{\watt}\right] \]
\textit{Exigence RE : IEE $\leq 0{,}23$ (circulateurs ErP 2015)}
\end{multicols}
\end{tcolorbox}

\finchapitre
